% ===========================================================================
% figures.tex -- shared TikZ content for all three candidate deck styles.
%
% Each style file defines these five colors before \input-ing this file, so
% the same drawings re-skin themselves per style:
%
%   fbCar     body fill for a freight car
%   fbLoco    body fill for a powered vehicle
%   fbAccent  the one saturated colour (forces, highlights)
%   fbLine    strokes, axes, rails
%   fbMuted   secondary text and de-emphasised strokes
% ===========================================================================

\usetikzlibrary{decorations.pathmorphing,arrows.meta,positioning,calc,patterns}

% --- a single vehicle body ------------------------------------------------
% #1 = x position, #2 = fill colour, #3 = label
\newcommand{\vehicle}[3]{%
  \draw[fill=#2,draw=fbLine,line width=0.7pt,rounded corners=1.5pt]
        (#1,0) rectangle ++(1.8,0.85);
  \node[font=\small,text=fbLine] at (#1+0.9,0.42) {#3};
  \draw[fill=fbLine,draw=fbLine] (#1+0.35,-0.02) circle (0.12);
  \draw[fill=fbLine,draw=fbLine] (#1+1.45,-0.02) circle (0.12);
}

% --- a coupler drawn as a coil spring ------------------------------------
% #1 = start x, #2 = end x
\newcommand{\coupler}[2]{%
  \draw[decorate,decoration={coil,aspect=0.55,segment length=2.2pt,amplitude=2.6pt},
        draw=fbAccent,line width=0.9pt] (#1,0.42) -- (#2,0.42);
}

% ===========================================================================
% \trainFBD -- free-body diagram of a three-vehicle consist on a grade.
% Used on slide 2 (plain) and slide 17 (symbol-annotated) of the plan.
% ===========================================================================
\newcommand{\trainFBD}[1][0]{%
\begin{tikzpicture}[scale=1.0,every node/.style={transform shape}]
  % --- inclined track -----------------------------------------------------
  \begin{scope}[rotate=6]
    \draw[fbLine,line width=1.1pt] (-0.4,-0.16) -- (9.4,-0.16);
    \foreach \x in {-0.2,0.25,...,9.2}
      \draw[fbMuted,line width=0.5pt] (\x,-0.16) -- ++(0,-0.16);

    \vehicle{0.0}{fbLoco}{Loco}
    \vehicle{3.2}{fbCar}{Car $i$}
    \vehicle{6.4}{fbCar}{Car $i{+}1$}

    \coupler{1.8}{3.2}
    \coupler{5.0}{6.4}

    % One force above each vehicle, so nothing collides:
    %   traction on the loco, resistance on car i, grade on car i+1.
    \draw[-{Stealth[length=6pt]},fbAccent,line width=1.5pt]
          (0.30,1.30) -- ++(1.20,0)
          node[midway,above,font=\small,text=fbAccent,inner sep=2.5pt]
          {$F_{\mathrm{trac}}$};
    \draw[{Stealth[length=6pt]}-,fbLine,line width=1.3pt]
          (3.50,1.30) -- ++(1.20,0)
          node[midway,above,font=\small,text=fbLine,inner sep=2.5pt]
          {$R(v)$};
    \draw[{Stealth[length=6pt]}-,fbMuted,line width=1.3pt]
          (6.70,1.30) -- ++(1.20,0)
          node[midway,above,font=\small,text=fbMuted,inner sep=2.5pt]
          {$m_i g\sin\theta$};

    % Coupler forces on the middle car, drawn BELOW the running gear so they
    % never collide with the coil springs. Opt in with \trainFBD[1].
    \ifnum#1=1
      \draw[fbMuted,line width=0.4pt,dotted] (3.2,0.36) -- (3.2,-0.78);
      \draw[fbMuted,line width=0.4pt,dotted] (5.0,0.36) -- (5.0,-0.78);
      \draw[-{Stealth[length=5pt]},fbAccent,line width=1.2pt] (3.2,-0.78) -- ++(0.75,0);
      \draw[-{Stealth[length=5pt]},fbAccent,line width=1.2pt] (5.0,-0.78) -- ++(-0.75,0);
      \node[font=\small,text=fbAccent,below,inner sep=2pt] at (3.55,-0.86) {$F_{i-1}$};
      \node[font=\small,text=fbAccent,below,inner sep=2pt] at (4.65,-0.86) {$F_{i}$};
    \fi
  \end{scope}

  % --- grade angle --------------------------------------------------------
  \draw[fbMuted,line width=0.6pt,dashed] (-0.45,-0.20) -- (2.6,-0.20);
  \draw[fbMuted,line width=0.6pt] (1.55,-0.20) arc (0:6:2.0);
  \node[font=\small,text=fbMuted] at (1.95,-0.04) {$\theta$};
\end{tikzpicture}}

% ===========================================================================
% \slackScene{n} -- one overlay of the slack run-in / run-out build.
%   n = 1  bunched (buff), couplers compressed, no force transmitted
%   n = 2  slack taken up, couplers at the edge of the deadband
%   n = 3  run-out (draft), force spike at take-up
% Used on slide 3 of the plan.
% ===========================================================================
\newcommand{\slackScene}[1]{%
\begin{tikzpicture}[scale=1.0,every node/.style={transform shape}]
  \def\spA{0.0}
  % Gaps stay wide enough that the coil still reads as a spring when bunched.
  \ifnum#1=1 \def\spB{2.30}\def\spC{4.60}\fi
  \ifnum#1=2 \def\spB{2.70}\def\spC{5.40}\fi
  \ifnum#1=3 \def\spB{3.10}\def\spC{6.20}\fi

  \draw[fbLine,line width=1.0pt] (-0.4,-0.16) -- (8.4,-0.16);
  \foreach \x in {-0.2,0.25,...,8.2}
    \draw[fbMuted,line width=0.5pt] (\x,-0.16) -- ++(0,-0.14);

  \vehicle{\spA}{fbLoco}{}
  \vehicle{\spB}{fbCar}{}
  \vehicle{\spC}{fbCar}{}

  \coupler{1.8}{\spB}
  \coupler{\spB+1.8}{\spC}

  % deadband bracket over the first coupler
  \draw[fbMuted,line width=0.7pt] (1.8,1.05) -- (1.8,1.55);
  \draw[fbMuted,line width=0.7pt] (\spB,1.05) -- (\spB,1.55);
  \draw[{Stealth[length=4pt]}-{Stealth[length=4pt]},fbMuted,line width=0.7pt]
        (1.8,1.42) -- (\spB,1.42);
  \node[font=\small,text=fbMuted,above,inner sep=2pt] at ({(1.8+\spB)/2},1.46)
        {Coupler Extension $\delta$};

  % --- state caption ------------------------------------------------------
  \ifnum#1=1
    \node[font=\normalsize\bfseries,text=fbMuted,anchor=west] at (0,-1.05)
      {Bunched: Slack Closed, No Force Transmitted};
  \fi
  \ifnum#1=2
    \node[font=\normalsize\bfseries,text=fbMuted,anchor=west] at (0,-1.05)
      {Slack Taken Up: Couplers Reach the End of the Slack};
  \fi
  \ifnum#1=3
    \node[font=\normalsize\bfseries,text=fbAccent,anchor=west] at (0,-1.05)
      {Run-Out: Force Spikes When the Slack Closes};
  \fi

  % --- coupler force trace, revealed progressively ------------------------
  \begin{scope}[xshift=0cm,yshift=-3.75cm]
    \draw[fbMuted,line width=0.7pt] (0,0) -- (7.6,0);
    \draw[fbMuted,line width=0.7pt] (0,-0.35) -- (0,1.75);
    \node[font=\small,text=fbMuted,rotate=90,anchor=south] at (-0.34,0.6)
      {Coupler Force $F_{i}$};
    \node[font=\small,text=fbMuted,anchor=north east] at (7.6,-0.08) {Time};
    \draw[fbMuted,line width=0.5pt,dashed] (0,0) -- (7.6,0);

    \ifnum#1>0
      \draw[fbLine,line width=1.2pt] (0,0) -- (2.4,0);
    \fi
    \ifnum#1>1
      \draw[fbLine,line width=1.2pt] (2.4,0) -- (3.9,0.04);
    \fi
    \ifnum#1>2
      \draw[fbAccent,line width=1.6pt]
        (3.9,0.04) .. controls (4.25,0.10) and (4.35,1.55) .. (4.6,1.55)
                   .. controls (4.9,1.55) and (5.0,0.35) .. (5.35,0.55)
                   .. controls (5.7,0.75) and (5.9,0.30) .. (6.3,0.48)
                   .. controls (6.8,0.62) and (7.1,0.40) .. (7.6,0.46);
      \draw[fbAccent,line width=0.6pt,dashed] (4.6,1.55) -- (0,1.55);
      \node[font=\small,text=fbAccent,anchor=south east,inner sep=2pt]
        at (7.6,1.60) {Peak Force: The Value the Safety Limit Applies To};
    \fi
  \end{scope}
\end{tikzpicture}}

% ===========================================================================
% \consistBand -- a long consist drawn as a wide, quiet band. Title slide.
% Deliberately low-contrast: it sits under the title block, not beside it.
% ===========================================================================
\newcommand{\consistBand}{%
\begin{tikzpicture}[scale=0.55,every node/.style={transform shape}]
  \draw[fbMuted!40,line width=0.8pt] (-0.3,-0.16) -- (16.9,-0.16);
  \foreach \x in {-0.1,0.4,...,16.7}
    \draw[fbMuted!30,line width=0.5pt] (\x,-0.16) -- ++(0,-0.12);
  \vehicle{0.0}{fbLoco}{}
  \foreach \i in {1,...,7}{%
    \pgfmathsetmacro{\px}{\i*2.05}
    \coupler{\px-0.25}{\px}
    \vehicle{\px}{fbCar}{}
  }
\end{tikzpicture}}

% ===========================================================================
% \surrogateFlow -- how a surrogate relates to the simulator it learns from.
% Slide 6. The point of the drawing is that physics stays the ground truth.
% ===========================================================================
\newcommand{\surrogateFlow}{%
\begin{tikzpicture}[
    node distance=0pt,
    box/.style={draw=uniNavy,line width=0.9pt,rounded corners=2pt,
                minimum width=3.5cm,minimum height=1.25cm,align=center,
                font=\small,text=uniInk},
    fill1/.style={box,fill=uniMist},
    fill2/.style={box,fill=chartBlue!18},
    arr/.style={-{Stealth[length=6pt]},uniNavy,line width=1.1pt}]

  \node[fill1] (sim) at (0,0) {High-Fidelity\\Physics Simulator};
  \node[fill1] (data) at (5.0,0) {Scenario\\Dataset};
  \node[fill2] (sur) at (10.0,0) {Learned\\Surrogate};

  \draw[arr] (sim) -- (data)
    node[midway,above,font=\scriptsize,text=uniSlate,inner sep=3pt] {Generates};
  \draw[arr] (data) -- (sur)
    node[midway,above,font=\scriptsize,text=uniSlate,inner sep=3pt] {Trains};

  % cost annotations beneath
  \node[font=\scriptsize,text=uniSlate,align=center] at (0,-1.15)
    {Hours of CPU Time\\per Scenario};
  \node[font=\scriptsize,text=chartRed,align=center] at (10.0,-1.15)
    {Milliseconds\\per Scenario};

  % the ground-truth loop back
  \draw[{Stealth[length=6pt]}-,uniCrimson,line width=1.0pt,dashed]
    (sim.north) .. controls +(0,1.15) and +(0,1.15) .. (sur.north);
  \node[font=\scriptsize,text=uniCrimson,align=center,fill=white,
        inner xsep=5pt,inner ysep=2pt] at (5.0,1.42)
    {Validated Against, Never Replaced By};
\end{tikzpicture}}

% ===========================================================================
% \actstrip{n} -- the roadmap / act-divider strip, act n highlighted.
% n = 0 shows the whole roadmap with nothing highlighted (slide 8).
% ===========================================================================
\newcommand{\actstrip}[1]{%
\begin{tikzpicture}[scale=1.0,every node/.style={transform shape}]
  \def\actlist{%
    1/Motivation, 2/Related Work, 3/Formulation, 4/Reference Simulator,
    5/Routes, 6/Data \& Tooling, 7/Architecture, 8/Training,
    9/Evaluation, 10/Control, 11/Conclusion}
  \foreach \n/\lbl [count=\i from 0] in {1/Motivation, 2/Related Work,
      3/Formulation, 4/Reference Simulator, 5/Routes, 6/Data \& Tooling,
      7/Architecture, 8/Training, 9/Evaluation, 10/Control, 11/Conclusion}{%
    \pgfmathsetmacro{\yy}{-\i*0.62}
    \ifnum#1=\n
      \node[anchor=west,font=\small\bfseries,text=white,fill=uniNavy,
            inner xsep=7pt,inner ysep=4pt,minimum width=5.4cm] at (0,\yy)
            {\makebox[4.9cm][l]{\n.~\lbl}};
    \else
      \node[anchor=west,font=\small,text=uniSlate,fill=uniMist!60,
            inner xsep=7pt,inner ysep=4pt,minimum width=5.4cm] at (0,\yy)
            {\makebox[4.9cm][l]{\n.~\lbl}};
    \fi
  }
\end{tikzpicture}}

% ===========================================================================
% \litMap{n} -- the four literatures this work sits between.
%   n = 0  all four quadrants plain
%   n = 1  the centre badge ("This Work") revealed
% Slide 9, re-shown highlighted on slide 15.
% ===========================================================================
\newcommand{\litMap}[1]{%
\begin{tikzpicture}[
    quad/.style={draw=uniNavy,line width=0.8pt,fill=uniMist!55,
                 minimum width=5.3cm,minimum height=2.1cm,align=center,
                 font=\small,text=uniInk,inner sep=4pt}]

  % NOTE: TikZ's \\ is not defined inside a nested {\scriptsize ...} group,
  % so the size change is applied inline rather than wrapped in braces.
  \node[quad] (a) at (-2.75, 1.25) {\textbf{Classical LTD Simulation}\\[0.15em]
        \scriptsize Physics-anchored, train-agnostic,\\ far too slow to optimize through};
  \node[quad] (b) at ( 2.75, 1.25) {\textbf{Data-Driven Rail Models}\\[0.15em]
        \scriptsize Fast, but tied to one consist\\ or one corridor};
  \node[quad] (c) at (-2.75,-1.25) {\textbf{Sequence \& Operator Models}\\[0.15em]
        \scriptsize General machinery, fixed token\\ or grid layout};
  \node[quad] (d) at ( 2.75,-1.25) {\textbf{Probabilistic Forecasting}\\[0.15em]
        \scriptsize Calibrated uncertainty,\\ no dynamics of its own};

  \ifnum#1=1
    \node[draw=uniCrimson,line width=1.4pt,fill=white,rounded corners=3pt,
          font=\small\bfseries,text=uniCrimson,inner xsep=9pt,inner ysep=6pt]
          at (0,0) {This Work};
  \fi
\end{tikzpicture}}

% ===========================================================================
% \mpStep{n} -- one message-passing step on a chain graph.
%   n = 1  the graph alone
%   n = 2  edge update from the two incident nodes
%   n = 3  node update aggregating incident edge messages
% Slide 13; reused for the architecture act.
% ===========================================================================
\newcommand{\mpStep}[1]{%
\begin{tikzpicture}[
    vtx/.style={circle,draw=uniNavy,line width=0.9pt,fill=uniMist,
                minimum size=0.78cm,font=\scriptsize,text=uniInk},
    vtxhot/.style={vtx,fill=chartBlue!30,draw=chartBlue,line width=1.2pt},
    edg/.style={uniNavy,line width=1.0pt},
    edghot/.style={chartRed,line width=1.6pt}]

  \foreach \i in {0,...,4}{%
    \pgfmathsetmacro{\px}{\i*2.3}
    \ifnum#1=3
      \ifnum\i=2 \node[vtxhot] (v\i) at (\px,0) {$h_\i$};
      \else      \node[vtx]    (v\i) at (\px,0) {$h_\i$};\fi
    \else
      \node[vtx] (v\i) at (\px,0) {$h_\i$};
    \fi
  }
  \foreach \i [evaluate=\i as \j using int(\i+1)] in {0,...,3}{%
    \ifnum#1>1
      \ifnum\i=1 \draw[edghot] (v\i) -- (v\j);
      \else\ifnum\i=2 \draw[edghot] (v\i) -- (v\j);
      \else \draw[edg] (v\i) -- (v\j);\fi\fi
    \else
      \draw[edg] (v\i) -- (v\j);
    \fi
  }
  % edge labels
  \foreach \i [evaluate=\i as \j using int(\i+1)] in {0,...,3}{%
    \pgfmathsetmacro{\px}{\i*2.3+1.15}
    \node[font=\scriptsize,text=uniSlate,above,inner sep=3pt] at (\px,0.10) {$e_{\i\j}$};
  }

  \ifnum#1=2
    \node[font=\small,text=chartRed,align=center] at (4.6,-1.35)
      {\textbf{Edge Update}\\[0.1em]
       {\scriptsize $e_{ij} \leftarrow \phi_e(e_{ij},\, h_i,\, h_j)$}};
  \fi
  \ifnum#1=3
    \node[font=\small,text=chartBlue!70!black,align=center] at (4.6,-1.35)
      {\textbf{Node Update}\\[0.1em]
       {\scriptsize $h_i \leftarrow \phi_h\bigl(h_i,\, \textstyle\sum_{j \in \mathcal{N}(i)} e_{ij}\bigr)$}};
  \fi
  \ifnum#1=1
    \node[font=\small,text=uniSlate,align=center] at (4.6,-1.35)
      {Vehicles Are Nodes, Couplers Are Edges\\[0.1em]
       {\scriptsize The Chain Is the Consist: Any Length, No Padding}};
  \fi
\end{tikzpicture}}

% ===========================================================================
% ACT III figures
% ===========================================================================

% --- \taskMap : three structured inputs -> operator -> four outputs -------
% Slide 16.
\newcommand{\taskMap}{%
\begin{tikzpicture}[
    inp/.style={draw=uniNavy,line width=0.8pt,fill=uniMist,rounded corners=2pt,
                minimum width=3.0cm,minimum height=0.95cm,align=center,
                font=\small,text=uniInk},
    outbox/.style={draw=uniSlate,line width=0.7pt,fill=white,rounded corners=2pt,
                minimum width=3.0cm,minimum height=0.8cm,align=center,
                font=\scriptsize,text=uniInk},
    arr/.style={-{Stealth[length=5pt]},uniNavy,line width=1.0pt}]

  \node[inp] (r) at (0, 1.45) {Route Profile $\mathcal{R}$};
  \node[inp] (c) at (0, 0.00) {Consist $\mathcal{C}$};
  \node[inp] (u) at (0,-1.45) {Control Plan $\mathcal{U}$};

  \node[draw=uniCrimson,line width=1.3pt,fill=white,rounded corners=3pt,
        minimum width=2.4cm,minimum height=3.4cm,align=center,
        font=\small\bfseries,text=uniCrimson] (op) at (4.5,0) {$F$};

  \node[outbox] (o1) at (9.1, 1.65) {Motion per Vehicle};
  \node[outbox] (o2) at (9.1, 0.55) {Force per Coupler};
  \node[outbox] (o3) at (9.1,-0.55) {Energy and Trip Time};
  \node[outbox] (o4) at (9.1,-1.65) {Safety Margin};

  \foreach \n in {r,c,u} \draw[arr] (\n) -- (op);
  \foreach \n in {o1,o2,o3,o4} \draw[arr] (op) -- (\n);
\end{tikzpicture}}

% --- \padWaste{n} : why a fixed token budget is the wrong bias ------------
% n = 1 the sequence model, n = 2 the graph model. Slide 18.
\newcommand{\padWaste}[1]{%
\begin{tikzpicture}[scale=1.0,every node/.style={transform shape}]
  \ifnum#1=1
    \node[font=\small\bfseries,text=uniInk,anchor=west] at (0,2.35)
      {Sequence Model: Fixed Input Length};
    % 150-slot strip, 11 real tokens
    \foreach \i in {0,...,29}{%
      \pgfmathsetmacro{\px}{\i*0.33}
      \ifnum\i<3
        \fill[chartBlue] (\px,1.35) rectangle ++(0.29,0.42);
      \else
        \fill[uniMist] (\px,1.35) rectangle ++(0.29,0.42);
        \draw[uniSlate!45,line width=0.3pt] (\px,1.35) rectangle ++(0.29,0.42);
      \fi
    }
    \node[font=\scriptsize,text=uniSlate,anchor=west] at (0,1.02)
      {11-car train: 3 tokens carry data, 27 slots are padding};
    \foreach \i in {0,...,29}{%
      \pgfmathsetmacro{\px}{\i*0.33}
      \fill[chartBlue] (\px,0.20) rectangle ++(0.29,0.42);
    }
    \node[font=\scriptsize,text=uniSlate,anchor=west] at (0,-0.13)
      {150-car train: the budget is exactly full};
    \node[font=\small,text=chartRed,anchor=west,align=left] at (0,-0.85)
      {The budget must be set for the longest train,\\
       so every short one wastes most of the tensor};
  \fi
  \ifnum#1=2
    \node[font=\small\bfseries,text=uniInk,anchor=west] at (0,2.35)
      {Graph Model: The Chain Is the Consist};
    \foreach \i in {0,...,3}{%
      \pgfmathsetmacro{\px}{\i*1.05}
      \node[circle,draw=uniNavy,line width=0.8pt,fill=uniMist,
            minimum size=0.5cm,font=\tiny] (a\i) at (\px,1.55) {};
    }
    \foreach \i [evaluate=\i as \j using int(\i+1)] in {0,...,2}
      \draw[chartBlue,line width=1.1pt] (a\i) -- (a\j);
    \node[font=\scriptsize,text=uniSlate,anchor=west] at (4.6,1.55)
      {short consist: 4 nodes, 3 edges};

    \foreach \i in {0,...,9}{%
      \pgfmathsetmacro{\px}{\i*1.05}
      \node[circle,draw=uniNavy,line width=0.8pt,fill=uniMist,
            minimum size=0.5cm,font=\tiny] (b\i) at (\px,0.30) {};
    }
    \foreach \i [evaluate=\i as \j using int(\i+1)] in {0,...,8}
      \draw[chartBlue,line width=1.1pt] (b\i) -- (b\j);
    \node[font=\scriptsize,text=uniSlate,anchor=west] at (10.0,0.30)
      {long consist: same weights};

    \node[font=\small,text=chartBlue!70!black,anchor=west,align=left] at (0,-0.85)
      {No padding, no masking. The parameter count\\
       does not depend on $N$};
  \fi
\end{tikzpicture}}

% --- \nodeTensor : the node (vehicle) feature matrix H -------------------
% Slide 19. Channels follow DATA_SCHEMA.md exactly.
\newcommand{\nodeTensor}{%
\begin{tikzpicture}[scale=1.0,every node/.style={transform shape}]
  \def\cw{0.62} \def\rh{0.46}
  % column headers
  \foreach \c/\lbl [count=\i from 0] in {%
      0/$x$, 1/$v$, 2/$z_{\mathrm{brk}}$, 3/$z_{\mathrm{trac}}$,
      4/Mass, 5/Davis $A$, 6/Davis $B$, 7/Davis $C$,
      8/Traction?, 9/$F^{\max}_{\mathrm{trac}}$, 10/$F^{\max}_{\mathrm{brk}}$}{%
    \pgfmathsetmacro{\px}{\i*\cw}
    \node[font=\scriptsize,text=uniSlate,rotate=52,anchor=west]
      at (\px+0.30,0.14) {\lbl};
  }
  % cells
  \foreach \r in {0,...,4}{%
    \pgfmathsetmacro{\py}{-\r*\rh}
    \foreach \i in {0,...,10}{%
      \pgfmathsetmacro{\px}{\i*\cw}
      \ifnum\i<4
        \fill[chartBlue!28] (\px,\py) rectangle ++(\cw-0.05,-\rh+0.05);
      \else
        \fill[uniMist] (\px,\py) rectangle ++(\cw-0.05,-\rh+0.05);
      \fi
      \draw[uniSlate!35,line width=0.3pt] (\px,\py) rectangle ++(\cw-0.05,-\rh+0.05);
    }
    \node[font=\scriptsize,text=uniSlate,anchor=east] at (-0.12,\py-\rh/2+0.03)
      {Vehicle \r};
  }
  \node[font=\scriptsize,text=uniSlate,anchor=east] at (-0.12,-5*\rh-0.06) {$\vdots$};
  % group rules -- a brace over this span renders almost flat
  \draw[chartBlue,line width=2.2pt] (0,-5*\rh-0.34) -- (4*\cw-0.05,-5*\rh-0.34);
  \node[font=\scriptsize,text=chartBlue!70!black,anchor=north]
    at (2*\cw,-5*\rh-0.44) {4 Dynamic};
  \draw[uniSlate,line width=2.2pt] (4*\cw,-5*\rh-0.34) -- (11*\cw-0.05,-5*\rh-0.34);
  \node[font=\scriptsize,text=uniSlate,anchor=north]
    at (7.5*\cw,-5*\rh-0.44) {7 Static: Stored Once, Not per Time Step};
\end{tikzpicture}}

% --- \edgeTensor : the edge (coupler) feature matrix E ------------------
% Slide 20.
\newcommand{\edgeTensor}{%
\begin{tikzpicture}[scale=1.0,every node/.style={transform shape}]
  \def\cw{0.62} \def\rh{0.46}
  \foreach \c/\lbl [count=\i from 0] in {%
      0/$\delta$, 1/$\dot\delta$, 2/$F$,
      3/$L_0$, 4/Slack, 5/$k_{\mathrm{draft}}$, 6/$c_{\mathrm{draft}}$,
      7/$k_{\mathrm{buff}}$, 8/$c_{\mathrm{buff}}$}{%
    \pgfmathsetmacro{\px}{\i*\cw}
    \node[font=\scriptsize,text=uniSlate,rotate=52,anchor=west]
      at (\px+0.30,0.14) {\lbl};
  }
  \foreach \r in {0,...,3}{%
    \pgfmathsetmacro{\py}{-\r*\rh}
    \foreach \i in {0,...,8}{%
      \pgfmathsetmacro{\px}{\i*\cw}
      \ifnum\i<3
        \fill[chartRed!25] (\px,\py) rectangle ++(\cw-0.05,-\rh+0.05);
      \else
        \fill[uniMist] (\px,\py) rectangle ++(\cw-0.05,-\rh+0.05);
      \fi
      \draw[uniSlate!35,line width=0.3pt] (\px,\py) rectangle ++(\cw-0.05,-\rh+0.05);
    }
    \node[font=\scriptsize,text=uniSlate,anchor=east] at (-0.12,\py-\rh/2+0.03)
      {Coupler \r};
  }
  \node[font=\scriptsize,text=uniSlate,anchor=east] at (-0.12,-4*\rh-0.06) {$\vdots$};
  \draw[chartRed,line width=2.2pt] (0,-4*\rh-0.34) -- (3*\cw-0.05,-4*\rh-0.34);
  \node[font=\scriptsize,text=chartRed,anchor=north]
    at (1.5*\cw,-4*\rh-0.44) {3 Dynamic};
  \draw[uniSlate,line width=2.2pt] (3*\cw,-4*\rh-0.34) -- (9*\cw-0.05,-4*\rh-0.34);
  \node[font=\scriptsize,text=uniSlate,anchor=north]
    at (6*\cw,-4*\rh-0.44) {6 Static};
\end{tikzpicture}}

% --- \imposedLearned : what the network outputs, and what is imposed ----
% Slides 23 and 67. The central design figure of the revised formulation.
\newcommand{\imposedLearned}[1]{%
\begin{tikzpicture}[
    ch/.style={draw=uniSlate,line width=0.7pt,rounded corners=2pt,
               minimum width=3.6cm,minimum height=0.85cm,align=center,
               font=\small,text=uniInk},
    learned/.style={ch,fill=chartBlue!25,draw=chartBlue,line width=1.1pt},
    imposed/.style={ch,fill=uniMist}]

  \node[font=\small\bfseries,text=uniInk] at (0,1.55) {State Channel};
  \node[font=\small\bfseries,text=uniInk] at (5.6,1.55) {Where It Comes From};

  \ifnum#1=1
    \node[ch] (x) at (0, 0.55) {$x$: position};
    \node[ch] (v) at (0,-0.45) {$v$: velocity};
    \node[ch] (zb) at (0,-1.45) {$z_{\mathrm{brk}},\, z_{\mathrm{trac}}$};
  \fi
  \ifnum#1>1
    \node[ch] (x) at (0, 0.55) {$x$: position};
    \node[ch] (v) at (0,-0.45) {$v$: velocity};
    \node[ch] (zb) at (0,-1.45) {$z_{\mathrm{brk}},\, z_{\mathrm{trac}}$};
    \node[imposed] (xr) at (5.6, 0.55) {Trapezoid $+$ Correction};
    \node[imposed] (zr) at (5.6,-1.45) {Exact Analytic Solution};
    \draw[-{Stealth[length=5pt]},uniSlate,line width=0.9pt] (x) -- (xr);
    \draw[-{Stealth[length=5pt]},uniSlate,line width=0.9pt] (zb) -- (zr);
  \fi
  \ifnum#1>2
    \node[learned] (vr) at (5.6,-0.45) {Learned Acceleration};
    \draw[-{Stealth[length=5pt]},chartBlue,line width=1.1pt] (v) -- (vr);
    \node[font=\small,text=chartBlue!70!black,align=center] at (2.8,-2.55)
      {The network outputs two scalars per vehicle,\\ not four channels};
  \fi
\end{tikzpicture}}

% --- \rolloutTensor : the T x N x d stack -------------------------------
% Slide 21.
\newcommand{\rolloutTensor}{%
\begin{tikzpicture}[scale=1.0,every node/.style={transform shape}]
  \foreach \s [count=\i from 0] in {0,1,2}{%
    \pgfmathsetmacro{\ox}{\i*0.42}
    \pgfmathsetmacro{\oy}{\i*0.30}
    \pgfmathsetmacro{\op}{100-\i*22}
    \fill[chartBlue!\op!white,opacity=0.95]
      (\ox,\oy) rectangle ++(3.3,2.0);
    \draw[uniNavy,line width=0.7pt] (\ox,\oy) rectangle ++(3.3,2.0);
  }
  \node[font=\scriptsize,text=uniSlate,anchor=north] at (1.65,-0.12) {$N$ Vehicles};
  \node[font=\scriptsize,text=uniSlate,rotate=90,anchor=south] at (-0.18,1.0) {$d$ Channels};
  \draw[-{Stealth[length=5pt]},uniSlate,line width=0.9pt]
    (3.55,0.15) -- ++(0.75,0.55);
  \node[font=\scriptsize,text=uniSlate,anchor=west] at (4.35,0.55) {$T$ Timesteps};
  \node[font=\small\bfseries,text=uniNavy,anchor=west] at (0.1,2.75)
    {$\mathbf{X} \in \mathbb{R}^{T \times N \times d}$};
\end{tikzpicture}}

% --- \consistOnGrade : one train spans several grade features ----------
% Slide 25. The point: route fields are sampled per vehicle, at x_i(t).
\newcommand{\consistOnGrade}{%
\begin{tikzpicture}[scale=1.0,every node/.style={transform shape}]
  % grade profile as a filled terrain curve
  \draw[uniSlate!35,line width=0.6pt] (0,-0.6) -- (13.0,-0.6);
  \draw[uniNavy,line width=1.3pt]
    (0,0.05) .. controls (2.2,0.05) and (2.8,0.95) .. (4.6,1.05)
             .. controls (6.4,1.15) and (6.8,0.35) .. (8.6,0.30)
             .. controls (10.4,0.25) and (11.2,1.00) .. (13.0,1.15);
  \fill[uniMist,opacity=0.65]
    (0,0.05) .. controls (2.2,0.05) and (2.8,0.95) .. (4.6,1.05)
             .. controls (6.4,1.15) and (6.8,0.35) .. (8.6,0.30)
             .. controls (10.4,0.25) and (11.2,1.00) .. (13.0,1.15)
             -- (13.0,-0.6) -- (0,-0.6) -- cycle;

  % vehicles riding it, each with its own local grade arrow
  \foreach \px/\py/\ang/\lbl in {%
      1.1/0.08/2/$\sin\theta_0$, 3.6/0.72/22/$\sin\theta_1$,
      6.1/1.05/-4/$\sin\theta_2$, 8.6/0.30/-14/$\sin\theta_3$,
      11.1/0.72/20/$\sin\theta_4$}{%
    \begin{scope}[shift={(\px,\py)},rotate=\ang]
      \draw[fill=fbCar,draw=uniNavy,line width=0.6pt,rounded corners=1pt]
        (-0.45,0.06) rectangle ++(0.9,0.42);
      \draw[fill=uniNavy,draw=uniNavy] (-0.25,0.04) circle (0.055);
      \draw[fill=uniNavy,draw=uniNavy] ( 0.25,0.04) circle (0.055);
    \end{scope}
    \node[font=\scriptsize,text=chartRed,anchor=south] at (\px,\py+0.60) {\lbl};
  }
  \node[font=\small,text=uniInk,anchor=west,align=left] at (0,-1.25)
    {Every vehicle samples the route at its own chainage $x_i(t)$.\\
     A 2\,km consist is on several grade features at once};
\end{tikzpicture}}

% --- \pipelineFig : the whole thesis as one flow -----------------------
% Slide 27, reused as the Act X divider and the closing slide.
\newcommand{\pipelineFig}[1]{%
\begin{tikzpicture}[
    src/.style={draw=uniSlate,line width=0.7pt,fill=white,rounded corners=2pt,
                minimum width=2.7cm,minimum height=0.78cm,align=center,
                font=\scriptsize,text=uniInk},
    stg/.style={draw=uniNavy,line width=0.9pt,fill=uniMist,rounded corners=2pt,
                minimum width=2.6cm,minimum height=1.15cm,align=center,
                font=\small,text=uniInk},
    hot/.style={stg,fill=chartBlue!22,draw=chartBlue,line width=1.3pt},
    arr/.style={-{Stealth[length=5pt]},uniNavy,line width=1.0pt}]

  \node[src] (ra) at (0, 1.25) {Route Acquisition};
  \node[src] (cl) at (0, 0.00) {Consist Library};
  \node[src] (cp) at (0,-1.25) {Control Profiles};

  \node[stg] (sim) at (4.1,0) {Reference\\Simulator};
  \node[stg] (cor) at (7.7,0) {Scenario\\Corpus};
  \ifnum#1=7 \node[hot] (sur) at (11.3,0) {Graph Network\\Simulator};
  \else      \node[stg] (sur) at (11.3,0) {Graph Network\\Simulator};\fi
  \ifnum#1=10 \node[hot] (mpc) at (14.9,0) {Fuel-Optimal\\MPC};
  \else       \node[stg] (mpc) at (14.9,0) {Fuel-Optimal\\MPC};\fi

  \foreach \n in {ra,cl,cp} \draw[arr] (\n) -- (sim);
  \draw[arr] (sim) -- (cor);
  \draw[arr] (cor) -- (sur);
  \draw[arr] (sur) -- (mpc);

  \node[font=\scriptsize,text=uniSlate,anchor=north] at (5.9,-0.72) {10,000 scenarios};
  \node[font=\scriptsize,text=chartRed,anchor=north] at (13.1,-0.72) {analytic gradients};
\end{tikzpicture}}

% ===========================================================================
% ACT IV figures
% ===========================================================================

% --- \stageLadder{n} : the eight validation stages, rung n highlighted ---
% Slide 28 shows the whole ladder (n = 0); slides 29-37 highlight one rung.
\newcommand{\stageLadder}[1]{%
\begin{tikzpicture}[scale=1.0,every node/.style={transform shape}]
  \foreach \n/\lbl [count=\i from 0] in {%
      1/{Single Powered Vehicle},
      2/{Two Bodies, Linear Coupler},
      3/{Nonlinear Coupler with Slack},
      4/{Three Bodies, Force Propagation},
      5/{General $N$-Vehicle Assembly},
      6/{Per-Vehicle Grade and Curvature},
      7/{Actuator Lag, Power-Limited Traction},
      8/{Scope vs Industrial High Fidelity}}{%
    \pgfmathsetmacro{\yy}{-\i*0.66}
    \ifnum#1=\n
      \node[anchor=west,font=\small\bfseries,text=white,fill=uniNavy,
            inner xsep=8pt,inner ysep=5pt,minimum width=6.6cm] at (0,\yy)
            {\makebox[6.0cm][l]{Stage \n~~\lbl}};
    \else
      \node[anchor=west,font=\small,text=uniSlate,fill=uniMist!60,
            inner xsep=8pt,inner ysep=5pt,minimum width=6.6cm] at (0,\yy)
            {\makebox[6.0cm][l]{Stage \n~~\lbl}};
    \fi
  }
\end{tikzpicture}}

% --- \scopeTable : what the simulator does and does not model -----------
% Slide 39, returned to on slide 94.
\newcommand{\scopeTable}{%
\begin{tikzpicture}[
    hdr/.style={font=\small\bfseries,text=uniInk},
    itm/.style={font=\footnotesize,text=uniInk,anchor=west}]
  \fill[chartBlue!12] (-0.25,0.45) rectangle (6.1,-3.95);
  \fill[uniMist!70]   (6.5,0.45) rectangle (12.9,-3.95);
  \node[hdr,anchor=west,text=chartBlue!70!black] at (-0.1,0.15) {In Scope};
  \node[hdr,anchor=west,text=uniSlate]           at (6.65,0.15) {Out of Scope};
  \foreach \t [count=\i from 0] in {%
      Longitudinal motion per vehicle,
      Nonlinear couplers with slack,
      Asymmetric buff and draft stiffness,
      Per-vehicle grade and curvature,
      First-order actuator lag,
      Power-limited traction}{%
    \pgfmathsetmacro{\yy}{-0.35-\i*0.58}
    \node[itm] at (-0.1,\yy) {\textbullet~~\t};
  }
  \foreach \t [count=\i from 0] in {%
      Lateral and vertical dynamics,
      Wheel--rail contact mechanics,
      Draft-gear internal states,
      Air-brake pipe propagation,
      Hysteresis and stick--slip,
      Calibration to real rolling stock}{%
    \pgfmathsetmacro{\yy}{-0.35-\i*0.58}
    \node[itm,text=uniSlate] at (6.65,\yy) {\textbullet~~\t};
  }
\end{tikzpicture}}

% --- \stagepips{n} : compact horizontal stage indicator ----------------
% The vertical \stageLadder is too tall to sit beside the stage figures,
% which are three-panel stacks. This is the in-slide version.
\newcommand{\stagepips}[1]{%
\begin{tikzpicture}[baseline]
  \foreach \i in {1,...,8}{%
    \pgfmathsetmacro{\px}{(\i-1)*0.52}
    \ifnum#1=\i
      \fill[uniNavy] (\px,0) rectangle ++(0.42,0.26);
      \node[font=\tiny,text=white] at (\px+0.21,0.13) {\i};
    \else
      \fill[uniMist] (\px,0) rectangle ++(0.42,0.26);
      \node[font=\tiny,text=uniSlate] at (\px+0.21,0.13) {\i};
    \fi
  }
\end{tikzpicture}}

% ===========================================================================
% ACT V figures
% ===========================================================================

% --- \routePipeline : acquisition -> resample -> elevation -> fields ----
% Slide 42. Fills the TODO placeholder at 05_route_acquisition.tex:14.
\newcommand{\routePipeline}{%
\begin{tikzpicture}[
    stg/.style={draw=uniNavy,line width=0.9pt,fill=uniMist,rounded corners=2pt,
                minimum width=2.45cm,minimum height=1.25cm,align=center,
                font=\footnotesize,text=uniInk},
    arr/.style={-{Stealth[length=5pt]},uniNavy,line width=1.0pt}]

  \node[stg] (a) at (0.0,0)  {\textbf{1}\\Acquisition};
  \node[stg] (b) at (3.3,0)  {\textbf{1.5}\\Resample};
  \node[stg] (c) at (6.6,0)  {\textbf{2}\\Elevation};
  \node[stg] (d) at (9.9,0)  {\textbf{3}\\Derived Fields};
  \node[stg] (e) at (13.2,0) {\textbf{4}\\Export};

  \foreach \x/\y in {a/b,b/c,c/d,d/e} \draw[arr] (\x) -- (\y);

  \foreach \n/\t in {a/{Corridors from a\\rail network graph},
                     b/{Uniform 10\,m grid},
                     c/{DEM samples along\\the alignment},
                     d/{$\sin\theta$, $\kappa$, $v_{\max}$},
                     e/{Route tensor for\\the simulator}}{%
    \node[font=\scriptsize,text=uniSlate,align=center,anchor=north]
      at ([yshift=-0.75cm]\n.south) {\t};
  }
\end{tikzpicture}}

% ===========================================================================
% ACT VI figures
% ===========================================================================

% --- \raillabLayers : screens -> services -> backend -------------------
% Slide 49. Fills the TODO placeholder at 06_raillab.tex:24.
\newcommand{\raillabLayers}{%
\begin{tikzpicture}[
    lyr/.style={draw=uniNavy,line width=0.9pt,rounded corners=2pt,
                minimum width=9.4cm,minimum height=1.15cm,align=center,
                font=\small,text=uniInk},
    arr/.style={{Stealth[length=5pt]}-{Stealth[length=5pt]},uniSlate,
                line width=1.0pt}]
  \node[lyr,fill=uniMist]        (s) at (0, 2.05) {\textbf{Screens}\\
    {\scriptsize Route \textbullet\ Consist \textbullet\ Controls \textbullet\ Simulate \textbullet\ Results}};
  \node[lyr,fill=chartBlue!18]   (v) at (0, 0.00) {\textbf{Services}\\
    {\scriptsize The only modules that import the numerical backend}};
  \node[lyr,fill=white]          (b) at (0,-2.05) {\textbf{Backend}\\
    {\scriptsize \texttt{simulator2} \textbullet\ \texttt{route\_generator}: unchanged, run without the GUI}};
  \draw[arr] (s) -- (v) node[midway,right,font=\scriptsize,text=uniSlate,inner sep=6pt]
    {screens call services};
  \draw[arr] (v) -- (b) node[midway,right,font=\scriptsize,text=uniSlate,inner sep=6pt]
    {services return plain data};
\end{tikzpicture}}

% --- \diskLayout : what a built corpus looks like on disk --------------
% Slide 59.
\newcommand{\diskLayout}{%
\begin{tikzpicture}[
    every node/.style={font=\footnotesize\ttfamily,text=uniInk,anchor=west},
    note/.style={font=\scriptsize\rmfamily,text=uniSlate,anchor=west}]
  \node at (0,0.0)   {data/v2/};
  \node at (0.5,-0.5)  {build\_config.json};   \node[note] at (5.0,-0.5)  {sampling ranges, git commit, seed};
  \node at (0.5,-1.0)  {index.parquet};        \node[note] at (5.0,-1.0)  {10,000 rows $\times$ 43 columns};
  \node at (0.5,-1.5)  {norm\_stats.json};     \node[note] at (5.0,-1.5)  {train split only};
  \node at (0.5,-2.0)  {splits.json};          \node[note] at (5.0,-2.0)  {explicit scenario ids per split};
  \node at (0.5,-2.5)  {routes/};              \node[note] at (5.0,-2.5)  {written once per corridor, not per scenario};
  \node at (0.5,-3.0)  {scenarios/<route>/};   \node[note] at (5.0,-3.0)  {one compressed .npz per scenario};
  \draw[uniSlate!40,line width=0.5pt] (0.25,-0.25) -- (0.25,-3.0);
  \foreach \y in {-0.5,-1.0,-1.5,-2.0,-2.5,-3.0}
    \draw[uniSlate!40,line width=0.5pt] (0.25,\y) -- (0.45,\y);
\end{tikzpicture}}
